\section{Tutorial: From the Data to a List of Targets}

This tutorial follows one analysis from the raw table to a shortlist
of diversification targets for a region, and then tests whether the
method behind the shortlist predicts where regions actually went. Each
step shows the code, what it prints and how to read it. The data are
synthetic (\fn{make\_sample\_data}), so every number below is reproduced
by running \texttt{examples/tutorial.py}; with your own data, replace
Step 1 and keep the rest.

The example follows region \texttt{L045}, which has an average number of
specializations. Region and activity labels are arbitrary codes.

\subsection{Step 1: From the Table to the Matrix}

The library starts from a long table: one row per region and
activity, with a value (employment, exports, patents).
\fn{pivot\_to\_matrix} turns it into the regions $\times$ activities
matrix that every function takes.

\begin{lstlisting}
df = ec.make_sample_data(n_locs=60, n_acts=40, seed=7)
print(df.head(3))
print(len(df), "rows |", df["loc"].nunique(), "regions |",
      df["act"].nunique(), "activities")

mat = ec.pivot_to_matrix(df, "loc", "act", "val")
print("matrix:", mat.shape)
\end{lstlisting}
\begin{lstlisting}[language={}, numbers=none, backgroundcolor=\color{white}, frame=leftline, framerule=2pt, rulecolor=\color{noteborder}, xleftmargin=1.5em]
    loc   act        val
0  L001  A001   2.450677
1  L002  A001  20.941576
2  L003  A001   3.358559
1450 rows | 60 regions | 40 activities
matrix: (60, 40)
\end{lstlisting}

With your own file the first line becomes
\texttt{df = pd.read\_csv("my\_data.csv")} and the three column names
change. Two things are worth checking before going on, because the
pivot hides them: pairs of region and activity that appear twice are
summed, and pairs that are absent become zero. Negative or missing
values are rejected by the functions with an error that names the
problem.

\subsection{Step 2: Who Is Specialized in What}

Raw values mostly reflect size: a large region has more of
everything. The RCA (the location quotient, when the data are
employment) divides the share of the activity in the region by its share
in the whole economy, and a region counts as specialized where
$\text{RCA} \geq 1$. The binary matrix $M$ records those
specializations; its row sums are the diversity of each region and its
column sums the ubiquity of each activity.

\begin{lstlisting}
rca = ec.rca(mat)
M = ec.mcp(mat)                    # 1 where RCA >= 1
diversity = ec.diversity(mat)      # specializations per region
ubiquity = ec.ubiquity(mat)        # specialized regions per activity

region = "L045"
print(rca.loc[region].nlargest(3).round(2))
print("specializations of", region, "=", int(diversity[region]))
print("median diversity =", diversity.median(),
      "| median ubiquity =", ubiquity.median())
\end{lstlisting}
\begin{lstlisting}[language={}, numbers=none, backgroundcolor=\color{white}, frame=leftline, framerule=2pt, rulecolor=\color{noteborder}, xleftmargin=1.5em]
act
A019    7.97
A009    3.62
A004    3.25
Name: L045, dtype: float64
specializations of L045 = 11
median diversity = 11.0 | median ubiquity = 17.5
\end{lstlisting}

\texttt{L045} is specialized in 11 of the
40 activities, and the median region in 11. Its
strongest specialization has an RCA of 7.97: the
activity weighs that many times more in the region than in the economy
as a whole. The threshold of 1 is the \texttt{threshold} argument of
every function that binarizes, and \texttt{use\_rca=False} skips the RCA
when the matrix is already an index or a 0/1 table.

\subsection{Step 3: Complexity, and How Far to Trust It}

The ECI ranks regions by the complexity of what they are
specialized in, and the PCI ranks activities by the complexity of the
regions that have them. \fn{eci\_pci} returns the two; the default is
the eigenvector method, and \texttt{method=} switches to reflections,
Fitness or GENEPY in the same call. \fn{eci\_diagnostics} then says how
firm the ranking is for this particular matrix.

\begin{lstlisting}
eci, pci = ec.eci_pci(mat)
order = eci.sort_values(ascending=False)
print(order.round(2).iloc[[0, 1, -2, -1]])
print("ECI of", region, "=", round(eci[region], 2), "| rank",
      int(eci.rank(ascending=False)[region]), "of", len(eci))

diag = ec.eci_diagnostics(mat, n_boot=100, seed=1)
keys = ["gap", "corr_diversity", "rank_stability_mean", "sign_flips"]
print(diag[keys].round(3))
\end{lstlisting}
\begin{lstlisting}[language={}, numbers=none, backgroundcolor=\color{white}, frame=leftline, framerule=2pt, rulecolor=\color{noteborder}, xleftmargin=1.5em]
loc
L002    1.63
L005    1.54
L059   -1.94
L056   -2.38
Name: eci, dtype: float64
ECI of L045 = -0.22 | rank 35 of 60
gap                    0.115
corr_diversity         0.773
rank_stability_mean    0.954
sign_flips             0.000
Name: eci_diagnostics, dtype: float64
\end{lstlisting}

Both indices are z-scores of this matrix: \texttt{L045} has an ECI
of $-$0.22, rank 35 of
60. Because they are relative, the ECI of two years or of two
classifications cannot be compared in level, only in rank; a region or
activity returned as \texttt{NaN} was left out for having no
specialization at all. The diagnostics are reassuring here: the ECI has
a correlation of 0.77 with diversity, and when a
fifth of the activities is dropped at random the ranking keeps a rank
correlation of 0.95 with the original and
never comes out inverted. In small or sparse matrices these numbers are
often much worse. If the correlation with diversity is close to zero,
the sign of the index is arbitrary, and it should be checked against
income or wages before ``high ECI'' is read as ``complex''.

\subsection{Step 4: What Is Within Reach}

Two activities are close when the same regions tend to be
specialized in both. From the proximity matrix $\phi$, relatedness
density measures, for each region and activity, how much of what is
related to the activity the region already has. It is the library's
basic measure of how feasible an activity is for a region.
\fn{relatedness\_density\_external} keeps only the activities the region
does not have, which are its options.

\begin{lstlisting}
phi = ec.product_proximity(mat)                  # proximity matrix
density = ec.relatedness_density(mat, phi=phi)   # 0 to 100
options = ec.relatedness_density_external(mat, phi=phi)

nearest = options.loc[region].nlargest(4)
print(nearest.round(1))

best = nearest.index[0]
why = ec.density_contributions(mat, region, best, phi=phi)
print(why.head(4).round(1))
print("sum of all contributions =", round(why.sum(), 1))
\end{lstlisting}
\begin{lstlisting}[language={}, numbers=none, backgroundcolor=\color{white}, frame=leftline, framerule=2pt, rulecolor=\color{noteborder}, xleftmargin=1.5em]
act
A017    38.1
A001    37.8
A005    35.6
A013    35.2
Name: L045, dtype: float64
act
A004    4.8
A014    4.0
A023    4.0
A003    3.8
Name: contribution, dtype: float64
sum of all contributions = 38.1
\end{lstlisting}

The closest activity \texttt{L045} lacks is \texttt{A017}, with
a density of 38.1: weighting by proximity, the region
has that percentage of what is related to it. \fn{density\_contributions}
opens the number. It lists the activities of the region that bring
\texttt{A017} closer, and their contributions add up to the density.
This is the sentence that goes into a regional profile:
``\texttt{A017} is within reach because the region already has
\texttt{A004}, \texttt{A014} and
\texttt{A023}''.

\subsection{Step 5: What Is Worth Reaching}

Feasibility says what is within reach, not what is worth
reaching. The second axis is a choice of the analyst: complexity, wages,
growth, emissions. \fn{desirability\_index} standardizes each indicator
and takes their weighted mean, with the weights written down. The wage
series below is simulated and stands for an external source.

\begin{lstlisting}
rng = np.random.default_rng(7)
wage = 1000 * np.exp(0.3 * pci + rng.normal(0, 0.3, len(pci)))
indicators = pd.DataFrame({"pci": pci, "wage": wage})
weights = {"pci": 0.6, "wage": 0.4}

desirability = ec.desirability_index(indicators, weights)
robust = ec.weight_sensitivity(indicators, weights, top=10, seed=1)
print(robust.head(4).round(2))
\end{lstlisting}
\begin{lstlisting}[language={}, numbers=none, backgroundcolor=\color{white}, frame=leftline, framerule=2pt, rulecolor=\color{noteborder}, xleftmargin=1.5em]
      score  rank  rank_mean  rank_p05  rank_p95  top_share
act                                                        
A035   1.96     1       1.20       1.0       3.0        1.0
A038   1.61     2       3.42       2.0      10.0        1.0
A037   1.19     3       4.83       3.0       9.0        1.0
A040   1.17     4       5.20       4.0       7.0        1.0
\end{lstlisting}

The weights are a judgment, so the ranking should not hang on
them. \fn{weight\_sensitivity} redraws the weights a thousand times
around the chosen ones and records where each activity lands.
\texttt{A035} is first with the chosen weights and stays
among the ten best in 100\,\% of the
draws. When \texttt{rank\_p05} and \texttt{rank\_p95} are far apart, the
position of the activity depends on the weights, and the report should
say so. An indicator that counts against an activity, such as emissions,
takes a negative weight.

\subsection{Step 6: The Shortlist}

With the two axes, each option of the region falls in one of
four quadrants. The cuts are the medians of the region's own options, so
``feasible'' means feasible compared with what else this region could
do. \fn{opportunity\_alignment} summarizes the picture in one number:
the correlation between the two axes among the region's options.

\begin{lstlisting}
quadrants = ec.strategy_quadrants(options, desirability)
print(quadrants.loc[region].value_counts())

alignment = ec.opportunity_alignment(options, desirability)
print("alignment of", region, "=", round(alignment[region], 2))

shortlist = pd.DataFrame({
    "density": options.loc[region],
    "desirability": desirability,
    "quadrant": quadrants.loc[region],
}).dropna()
desirable = shortlist["quadrant"].isin(["short jump", "long bet"])
wanted = shortlist[desirable]
print(wanted.sort_values("density", ascending=False).head(4).round(2))
\end{lstlisting}
\begin{lstlisting}[language={}, numbers=none, backgroundcolor=\color{white}, frame=leftline, framerule=2pt, rulecolor=\color{noteborder}, xleftmargin=1.5em]
L045
low value     13
long bet      13
avoid          2
short jump     1
Name: count, dtype: int64
alignment of L045 = -0.87
      density  desirability    quadrant
act                                    
A024    30.24          0.91  short jump
A030    28.85          0.71    long bet
A026    28.40          0.52    long bet
A028    25.92          0.95    long bet
\end{lstlisting}

Of the 29 activities \texttt{L045} lacks, only
1 is a short jump (feasible and desirable);
13 are long bets
(desirable but far) and
13 are within reach
but of low value. The alignment of $-$0.87 says the same:
among this region's options, the closer an activity is, the less
desirable it tends to be. This is the usual situation of regions of low
or middling complexity, and it changes what the shortlist has to be.
With so few short jumps, the useful list is the desirable half of the
options ordered by density: the one short jump,
\texttt{A024}, followed by the long bets that are least
far, \texttt{A030} and \texttt{A026},
with densities between 28.4 and
30.2. In a region with positive alignment the
short jumps would fill the list by themselves. How many unrelated bets
to make, and in which order, is what \fn{optimize\_sequence} answers,
and \fn{eci\_optimization} returns the set of new activities that
reaches a target ECI with the least effort (see the sections on
Strategic Diffusion and ECI Optimization).

\subsection{Step 7: Reading the Result Through Policy Lists}

Policies are usually written as lists of activities: a mission,
a value chain, a green taxonomy. The list functions read the same matrix
through a list. Below, the last fifteen activities stand for a mission
and three groups of activities for the links of a chain.

\begin{lstlisting}
acts = list(mat.columns)
lists = {"mission": acts[25:]}
feasibility = density / 100            # shares, 0 to 1
tab = ec.list_complexity_potential(mat, pci, feasibility, lists)
print(tab.loc[region].round(2))

chain = {"inputs": acts[0:6], "components": acts[15:21],
         "assembly": acts[33:39]}
links = ec.missing_links(mat, chain, feasibility, cut=0.4)
print(links.loc[region])
\end{lstlisting}
\begin{lstlisting}[language={}, numbers=none, backgroundcolor=\color{white}, frame=leftline, framerule=2pt, rulecolor=\color{noteborder}, xleftmargin=1.5em]
         presence  potential  n_present  n_list
list                                           
mission       0.0       0.23          0      15
inputs        present
components    present
assembly      missing
Name: L045, dtype: object
\end{lstlisting}

\texttt{L045} is specialized in
0 of the
15 activities of the
mission, which here are the most complex ones. \texttt{presence} adds
the PCI of the activities it has from the list, and \texttt{potential}
is the mean PCI of those it lacks, each weighted by its density; with a list of green products these are the green
complexity index and potential. For the chain, a link is
\texttt{present} when the region is specialized in at least one of its
activities, \texttt{feasible} when some activity of the link has a
density of 40\,\% or more, and \texttt{missing} otherwise.
\fn{transition\_risk} applies the same logic to activities expected to
shrink.

\subsection{Step 8: Does Density Predict Entry?}

Everything so far assumes that density says where regions will
go. With two periods this can be tested: compute the scores with the
first period and see whether they rank the entries observed by the
second. Here the second period is simulated, with faster growth where
density is high; with real data, load the later year in its place.

\begin{lstlisting}
size = np.outer(mat.sum(axis=1), mat.sum(axis=0)) / mat.values.sum()
mat_t1 = (mat * rng.lognormal(0, 0.3, mat.shape)
          + size * rng.lognormal(0, 1.0, mat.shape)
          * 3 * (density / 100) ** 2)

entered, at_risk = ec.entry_targets(rca, ec.rca(mat_t1))
print(int(at_risk.values.sum()), "cells at risk |",
      int(entered.sum().sum()), "entries")

scores = {"density": density,
          "relative": ec.relative_relatedness(mat, phi=phi),
          "persistence": rca}
val = ec.validate_entry(scores, entered, k=(5,),
                        baseline="persistence", n_boot=200, seed=1)
cols = ["auc", "auc_within", "precision@5"]
print(val["value"].unstack()[cols].round(3))
print(val.loc[("density", "auc")].round(3))
\end{lstlisting}
\begin{lstlisting}[language={}, numbers=none, backgroundcolor=\color{white}, frame=leftline, framerule=2pt, rulecolor=\color{noteborder}, xleftmargin=1.5em]
1400 cells at risk | 67 entries
metric         auc  auc_within  precision@5
score                                      
density      0.788       0.646        0.073
persistence  0.686       0.626        0.097
relative     0.639       0.646        0.073
value      0.788
lo         0.734
hi         0.839
diff       0.102
diff_lo    0.039
diff_hi    0.170
Name: (density, auc), dtype: float64
\end{lstlisting}

A cell is at risk when its RCA is below 0.5 at the start, and it
enters when the RCA reaches 1; the gap avoids counting cells that only
hover around 1. Of 1400 cells at risk,
67 entered. The AUC is the probability that an
entry is ranked above a non-entry: 0.79
for density and 0.69 for
persistence, the initial RCA. \texttt{auc\_within} computes the same
inside each region and averages. It is the number that matters for
choosing among the options of one region, and it is lower because part
of the pooled AUC comes from telling regions apart: inside regions,
density reaches 0.65 and
persistence 0.63, a much
smaller distance.
\texttt{precision@5} is the share of each region's five best-ranked
options that entered. The last block compares density with persistence:
a gain of 0.10 in AUC, with a 95\,\%
interval from 0.04 to
0.17 obtained by resampling
regions. For the test to be honest, every score has to use data from the
first period only, the proximity matrix included.

\subsection{Step 9: A Prediction Model and Expected Jobs (Optional)}

With several measures of feasibility, a model can combine them.
\fn{build\_features} stacks the layers into one table of cells and adds
controls; \fn{predict\_entry} cross-validates by region, so that the
probability of each cell comes from a model that never saw its region.
This step needs scikit-learn
(\texttt{pip install "econcomplex[ml]"}).

\begin{lstlisting}
layers = {"density": density, "implied": ec.implied_rca(mat)}
X = ec.build_features(mat, layers, at_risk=at_risk)
scores["model"] = ec.predict_entry(X, entered, model="logit", seed=1)
out = ec.validate_entry(scores, entered, k=(5,))["value"].unstack()
print(out.loc[["density", "model"], ["auc", "auc_within"]].round(3))

importance = ec.explain_entry(X, entered, model="logit", seed=1)
print(importance.head(3).round(3))

# next window: fit on t0 -> t1, predict from what is observed at t1
phi_t1 = ec.product_proximity(mat_t1)
layers_t1 = {"density": ec.relatedness_density(mat_t1, phi=phi_t1),
             "implied": ec.implied_rca(mat_t1)}
X_t1 = ec.build_features(mat_t1, layers_t1,
                         at_risk=ec.rca(mat_t1) < 0.5)
prob = ec.predict_entry(X, entered, model="logit",
                        new_features=X_t1, calibrate=True)
jobs = ec.potential_employment(prob, mat_t1)
print(jobs.groupby("region")["expected"].sum().nlargest(3).round(1))
\end{lstlisting}
\begin{lstlisting}[language={}, numbers=none, backgroundcolor=\color{white}, frame=leftline, framerule=2pt, rulecolor=\color{noteborder}, xleftmargin=1.5em]
metric     auc  auc_within
score                     
density  0.788       0.646
model    0.792       0.643
                 importance    std  p_value  significant
density               0.273  0.099    0.020         True
log_region_size       0.040  0.037    0.137        False
rca                   0.023  0.021    0.137        False
region
L019    540.4
L016    530.3
L017    519.1
Name: expected, dtype: float64
\end{lstlisting}

The model enters \fn{validate\_entry} as one more score. Here it
ties with plain density (an AUC of 0.792 against
0.788, and
0.643 against
0.646 inside regions), which is what
one expects from data that were simulated from density alone. A model is worth its extra machinery
only when it wins this comparison out of sample. \fn{explain\_entry}
shuffles one feature at a time in the held-out regions and reports the
fall in AUC: density carries the prediction
(0.27) and no other feature is
statistically distinguishable from zero. For the next window the same function takes
\texttt{new\_features}: the model is fitted on the first window and
applied to what is observed at the end of it. \fn{potential\_employment}
then multiplies each calibrated probability by the size the activity
would have in the region with a location quotient of 1, and the sum by
region is the expected employment in new specializations.

\subsection{Before Using the Results}

A few points account for most of the mistakes in applied work with these
indicators.

\begin{itemize}[leftmargin=1.5em]
  \item ECI and PCI are relative to one matrix. To follow a region over
    time, compare ranks, or project a fixed PCI on each year with
    \fn{subnational\_eci}.
  \item \texttt{NaN} means that the unit was left out (no
    specialization, or not an option); it is not a zero.
  \item In sparse data, where many regions have one or two
    specializations, compute complexity on the connected core with
    \texttt{dmin=2, umin=2}.
  \item Density comes in percent (0 to 100) by default and as a share
    with \texttt{percent=False}. Cuts such as the one of
    \fn{missing\_links} are read on the scale of the matrix you pass.
  \item External vectors (GDP, population, wages, PCI from another
    source) are matched to the matrix by label, in any order; a label
    that is missing raises an error instead of being filled.
  \item In any test of prediction, the proximity matrix and every score
    use the first period only.
  \item State the weights of the desirability index and report how the
    ranking moves when they change.
\end{itemize}
